
est constante sur @@M0@@. Comme @@M1@@ est projectif, il existe une immersion fermée @@M2@@, où @@M3@@ est un fibré projectif sur @@M4@@, et comme @@M5@@ a un faisceau ample, on peut supposer que @@M6@@. On a donc une factorisation :
@@M7@@
Comme @@M8@@ est d'intersection complète, l'immersion @@M9@@ est régulière (1.2).

D'après 1.7, @@M10@@ est un morphisme parfait. De plus, @@M11@@ a un faisceau ample et par suite @@M12@@ et @@M13@@ également. On peut alors définir un homomorphisme additif :
@@M14@@
(IV 2.12), les @@M15@@ définis en termes de faisceaux localement libres de type fini, ou de complexes parfaits, étant identiques. On peut également définir des homomorphismes additifs :
@@M16@@
et on a : @@M17@@.

Soit @@M18@@ la dimension relative virtuelle de @@M19@@ sur @@M20@@. La codimension de @@M21@@ est alors @@M22@@. D'après VII 4.6, l'homomorphisme @@M23@@ donne après tensorisation par @@M24@@ des inclusions :
@@M25@@
D'après VI 5.8, l'homomorphisme @@M26@@ donne des inclusions :
@@M27@@
En composant les deux, on obtient donc:

